\documentclass[11pt]{article}

\usepackage{geometry}
\usepackage{amsmath}
\usepackage{booktabs}
\usepackage{siunitx}
\usepackage{microtype}
\usepackage{url}
\geometry{margin=1in}

\DeclareSIUnit\calorie{cal}
\DeclareSIUnit\kcal{\kilo\calorie}
\DeclareSIUnit\atm{atm}

\title{Room Heating from the Cohesive Energy of Rocksalt}
\author{Yao Luo}
\date{26 July 2026}

\begin{document}
\maketitle

\section{Problem}

\begin{quote}
\itshape
Imagine yourself to be in a room, say a typical living room or class room. You
get a box with free sodium atoms (Na) and another box with free chlorine atoms
(Cl). You connect the two boxes, and let the atoms react with each other. The
result is one teaspoon of rocksalt. How much has the temperature of the room
increased? Assume the cohesive energy of rocksalt to be 187 kcal/mol.
\end{quote}

\section{Model and assumptions}

\begin{enumerate}
  \item The boxes are thin, indestructible, perfectly conducting, and of
        negligible heat capacity and volume. All energy released by the reaction is therefore transferred to the room environment.
  \item The room is sealed and perfectly insulated, so the air is heated at
        constant volume and none of the heat escapes to the walls.
  \item The room's dimensions are approximated as that of a typical university
        classroom (\qty{50}{\meter\squared} floor area, \qty{3}{\meter} high).
  \item The air is an ideal diatomic gas, with $C_V = \tfrac{5}{2}R
    = \qty{20.8}{\joule\per\mole\per\kelvin}$.
  \item The reactants begin as free neutral atoms and end as crystalline NaCl, with the reaction going to completion, so the full cohesive energy is released as heat.
  \item The heat capacity of the salt itself is negligible against that of the air.
\end{enumerate}

\begin{table}[htbp]
  \centering
  \caption{Input quantities.}\label{tab:givens}
  \begin{tabular}{llS[table-format=3.2]ll}
    \toprule
    Quantity & Symbol & {Value} & Unit & Source \\
    \midrule
    Room volume            & $V$          & 150   & \unit{\meter\cubed} & [assumed] \\
    Ambient temperature    & $T_0$        & 300   & \unit{\kelvin} & [assumed] \\
    Ambient pressure       & $P$          & 1     & \unit{\atm} & [assumed] \\
    Molar heat capacity of air & $C_V$    & 20.8  & \unit{\joule\per\mole\per\kelvin} & \cite{feynman} \\
    \midrule
    Volume of one teaspoon & $V_\text{s}$ & 4.9   & \unit{\milli\litre} & \cite{hb44} \\
    Density of NaCl        & $\rho$       & 2.17  & \unit{\gram\per\cubic\cm} & \cite{halite} \\
    Molar mass of NaCl     & $M$          & 58.44 & \unit{\gram\per\mole} & \cite{ciaaw} \\
    Cohesive energy of NaCl & $E_\text{coh}$ & 187 & \unit{\kcal\per\mole} & [given] \\
    \bottomrule
  \end{tabular}
\end{table}

\section{Solution}

\subsection{Amount of rocksalt formed}

One teaspoon is $V_\text{s} = \qty{4.9}{\milli\litre} = \qty{4.9}{\cubic\cm}$, so
\begin{equation}
  n_\text{NaCl} = \frac{\rho V_\text{s}}{M}
  = \frac{\qty{2.17}{\gram\per\cubic\cm} \times \qty{4.9}{\cubic\cm}}
         {\qty{58.44}{\gram\per\mole}}
  = \qty{0.182}{\mole}.
  \label{eq:moles}
\end{equation}

\subsection{Energy released}

The cohesive energy is the energy required to break the crystal into free
neutral atoms,
\begin{equation}
  E_\text{coh} = \left(E_\text{Na} + E_\text{Cl}\right) - E_\text{NaCl}
  = \qty{187}{\kcal\per\mole},
  \label{eq:ecoh}
\end{equation}
so exactly this much is liberated per mole when the crystal forms from those
same free atoms:
\begin{equation}
  Q = n_\text{NaCl}\,E_\text{coh}
  = \qty{0.182}{\mole} \times \qty{187}{\kcal\per\mole}
  = \qty{34.0}{\kcal} = \qty{142}{\kilo\joule}.
  \label{eq:Q}
\end{equation}

\subsection{Amount of air in the room}

With $V = \qty{50}{\meter\squared} \times \qty{3}{\meter} =
\qty{150}{\cubic\meter}$, the ideal gas law gives~\cite{codata}
\begin{equation}
  n_\text{air} = \frac{PV}{RT_0}
  = \frac{\qty{101325}{\pascal} \times \qty{150}{\cubic\meter}}
         {\qty{8.314}{\joule\per\mole\per\kelvin} \times \qty{300}{\kelvin}}
  = \qty{6.09e3}{\mole}.
  \label{eq:nair}
\end{equation}

\subsection{Temperature rise}

Heating that air at constant volume,
\begin{equation}
  \Delta T = \frac{Q}{n_\text{air} C_V}
  = \frac{\qty{1.424e5}{\joule}}
         {\qty{6.09e3}{\mole} \times \qty{20.8}{\joule\per\mole\per\kelvin}}
  \quad\Longrightarrow\quad
  \boxed{\;\Delta T \approx \qty{1.1}{\kelvin}\;}
  \label{eq:answer}
\end{equation}

\section{Discussion}

A teaspoon of salt is only $\sim\!\qty{0.18}{\mole}$, but the room holds
$\sim\!\qty{6.1e3}{\mole}$ of air, over four orders of magnitude more. The
$\qty{187}{\kcal\per\mole}$ of cohesive energy is enormous on an atomic scale
($\approx\qty{8}{\electronvolt}$ per formula unit). Despite the small quantity
of reactants compared to the room volume, the reaction is so exothermic that the
temperature change is still noticeable on the latter scale.

One modelling choice is significant to the magnitude of the answer: whether the
air is heated at constant volume or constant pressure. A sealed room fixes $V$,
so $C_V$ is the appropriate heat capacity. A room vented to the outside would
instead expand against the atmosphere and require $C_P = \tfrac{7}{2}R =
\qty{29.1}{\joule\per\mole\per\kelvin}$, lowering the answer to $\Delta T
\approx \qty{0.8}{\kelvin}$.

\clearpage
\begin{thebibliography}{9}
\raggedright

\bibitem{ciaaw}
IUPAC Commission on Isotopic Abundances and Atomic Weights,
``Abridged Standard Atomic Weights 2024.''
$A_\mathrm{r}^{\circ}(\mathrm{Na}) = 22.990$,
$A_\mathrm{r}^{\circ}(\mathrm{Cl}) = 35.45$.
\url{https://www.ciaaw.org/abridged-atomic-weights.htm}
(accessed 26 July 2026).

\bibitem{halite}
``Halite, NaCl,'' \emph{Handbook of Mineralogy},
Mineral Data Publishing (2001--2005).
$D(\text{meas.}) = \qty{2.168}{\gram\per\cubic\cm}$.
\url{https://www.handbookofmineralogy.org/pdfs/halite.pdf}
(accessed 26 July 2026).

\bibitem{hb44}
National Institute of Standards and Technology,
\emph{Specifications, Tolerances, and Other Technical Requirements for
Weighing and Measuring Devices}, NIST Handbook~44, 2026 edition,
Appendix~C. One measuring teaspoon is $1\tfrac{1}{3}$ fluid drams
(p.~C-26) and one fluid dram is \qty{3.696691}{\milli\litre} (p.~C-16),
giving \qty{4.9289}{\milli\litre}.
\url{https://doi.org/10.6028/NIST.HB.44-2026}
(accessed 26 July 2026).

\bibitem{codata}
National Institute of Standards and Technology,
``Fundamental Physical Constants --- Complete Listing,''
2022 CODATA adjustment.
$R = 8.314\,462\,618\ldots\ \mathrm{J\,mol^{-1}\,K^{-1}}$ and
standard atmosphere $= \qty{101325}{\pascal}$, both exact.
\url{https://physics.nist.gov/cuu/Constants/Table/allascii.txt}
(accessed 26 July 2026).

\bibitem{feynman}
R.~P. Feynman, R.~B. Leighton, and M.~Sands,
\emph{The Feynman Lectures on Physics}, New Millennium Edition,
Vol.~I, Ch.~40, \S40--5, ``The specific heats of gases,''
for the rigid-diatomic result $C_V = \tfrac{5}{2}R$.
\url{https://www.feynmanlectures.caltech.edu/I\_40.html}
(accessed 26 July 2026).

\end{thebibliography}

\end{document}
